% Options for packages loaded elsewhere
% Options for packages loaded elsewhere
\PassOptionsToPackage{unicode}{hyperref}
\PassOptionsToPackage{hyphens}{url}
\PassOptionsToPackage{dvipsnames,svgnames,x11names}{xcolor}
%
\documentclass[
  letterpaper,
  DIV=11,
  numbers=noendperiod]{scrartcl}
\usepackage{xcolor}
\usepackage{amsmath,amssymb}
\setcounter{secnumdepth}{-\maxdimen} % remove section numbering
\usepackage{iftex}
\ifPDFTeX
  \usepackage[T1]{fontenc}
  \usepackage[utf8]{inputenc}
  \usepackage{textcomp} % provide euro and other symbols
\else % if luatex or xetex
  \usepackage{unicode-math} % this also loads fontspec
  \defaultfontfeatures{Scale=MatchLowercase}
  \defaultfontfeatures[\rmfamily]{Ligatures=TeX,Scale=1}
\fi
\usepackage{lmodern}
\ifPDFTeX\else
  % xetex/luatex font selection
\fi
% Use upquote if available, for straight quotes in verbatim environments
\IfFileExists{upquote.sty}{\usepackage{upquote}}{}
\IfFileExists{microtype.sty}{% use microtype if available
  \usepackage[]{microtype}
  \UseMicrotypeSet[protrusion]{basicmath} % disable protrusion for tt fonts
}{}
\makeatletter
\@ifundefined{KOMAClassName}{% if non-KOMA class
  \IfFileExists{parskip.sty}{%
    \usepackage{parskip}
  }{% else
    \setlength{\parindent}{0pt}
    \setlength{\parskip}{6pt plus 2pt minus 1pt}}
}{% if KOMA class
  \KOMAoptions{parskip=half}}
\makeatother
% Make \paragraph and \subparagraph free-standing
\makeatletter
\ifx\paragraph\undefined\else
  \let\oldparagraph\paragraph
  \renewcommand{\paragraph}{
    \@ifstar
      \xxxParagraphStar
      \xxxParagraphNoStar
  }
  \newcommand{\xxxParagraphStar}[1]{\oldparagraph*{#1}\mbox{}}
  \newcommand{\xxxParagraphNoStar}[1]{\oldparagraph{#1}\mbox{}}
\fi
\ifx\subparagraph\undefined\else
  \let\oldsubparagraph\subparagraph
  \renewcommand{\subparagraph}{
    \@ifstar
      \xxxSubParagraphStar
      \xxxSubParagraphNoStar
  }
  \newcommand{\xxxSubParagraphStar}[1]{\oldsubparagraph*{#1}\mbox{}}
  \newcommand{\xxxSubParagraphNoStar}[1]{\oldsubparagraph{#1}\mbox{}}
\fi
\makeatother


\usepackage{longtable,booktabs,array}
\usepackage{calc} % for calculating minipage widths
% Correct order of tables after \paragraph or \subparagraph
\usepackage{etoolbox}
\makeatletter
\patchcmd\longtable{\par}{\if@noskipsec\mbox{}\fi\par}{}{}
\makeatother
% Allow footnotes in longtable head/foot
\IfFileExists{footnotehyper.sty}{\usepackage{footnotehyper}}{\usepackage{footnote}}
\makesavenoteenv{longtable}
\usepackage{graphicx}
\makeatletter
\newsavebox\pandoc@box
\newcommand*\pandocbounded[1]{% scales image to fit in text height/width
  \sbox\pandoc@box{#1}%
  \Gscale@div\@tempa{\textheight}{\dimexpr\ht\pandoc@box+\dp\pandoc@box\relax}%
  \Gscale@div\@tempb{\linewidth}{\wd\pandoc@box}%
  \ifdim\@tempb\p@<\@tempa\p@\let\@tempa\@tempb\fi% select the smaller of both
  \ifdim\@tempa\p@<\p@\scalebox{\@tempa}{\usebox\pandoc@box}%
  \else\usebox{\pandoc@box}%
  \fi%
}
% Set default figure placement to htbp
\def\fps@figure{htbp}
\makeatother





\setlength{\emergencystretch}{3em} % prevent overfull lines

\providecommand{\tightlist}{%
  \setlength{\itemsep}{0pt}\setlength{\parskip}{0pt}}



 


\KOMAoption{captions}{tableheading}
\makeatletter
\@ifpackageloaded{caption}{}{\usepackage{caption}}
\AtBeginDocument{%
\ifdefined\contentsname
  \renewcommand*\contentsname{Table of contents}
\else
  \newcommand\contentsname{Table of contents}
\fi
\ifdefined\listfigurename
  \renewcommand*\listfigurename{List of Figures}
\else
  \newcommand\listfigurename{List of Figures}
\fi
\ifdefined\listtablename
  \renewcommand*\listtablename{List of Tables}
\else
  \newcommand\listtablename{List of Tables}
\fi
\ifdefined\figurename
  \renewcommand*\figurename{Figure}
\else
  \newcommand\figurename{Figure}
\fi
\ifdefined\tablename
  \renewcommand*\tablename{Table}
\else
  \newcommand\tablename{Table}
\fi
}
\@ifpackageloaded{float}{}{\usepackage{float}}
\floatstyle{ruled}
\@ifundefined{c@chapter}{\newfloat{codelisting}{h}{lop}}{\newfloat{codelisting}{h}{lop}[chapter]}
\floatname{codelisting}{Listing}
\newcommand*\listoflistings{\listof{codelisting}{List of Listings}}
\makeatother
\makeatletter
\makeatother
\makeatletter
\@ifpackageloaded{caption}{}{\usepackage{caption}}
\@ifpackageloaded{subcaption}{}{\usepackage{subcaption}}
\makeatother
\usepackage{bookmark}
\IfFileExists{xurl.sty}{\usepackage{xurl}}{} % add URL line breaks if available
\urlstyle{same}
\hypersetup{
  colorlinks=true,
  linkcolor={blue},
  filecolor={Maroon},
  citecolor={Blue},
  urlcolor={Blue},
  pdfcreator={LaTeX via pandoc}}


\author{}
\date{}
\begin{document}


\section{On Analogies of Dynamical
Systems}\label{on-analogies-of-dynamical-systems}

At a fundamental level, no system is truly static and nothing happens
for no reason. The nature of a dynamical system is that it changes - for
some reason. The study of dynamical systems seeks to describe the
behavior of systems as they change. As one might imagine, there are a
multitude of factors and influences which affect the behavior of various
systems.Of specific interest is that which encourages or discourages a
system to change. What motivations or hesitations, attractions or
repulsions influence its behavior? Are these influences governed by the
nature of the system itself or the environment in which the system
evolves? Or both? If so, which has a dominant effect? This is a question
of ``how?'', not of ``why?''

It is not merely that change is interesting, but that it is essential to
life and survival - it is the driving force. When surveying for a target
or prey, it is optimal for a predator to scan the scene with their eyes
rather than staring in one direction. Focusing attention in one location
creates a central bias which tends to ignore the periphery - where
change is more likely to occur. Scanning attention across locations
creates a peripheral bias where change is more likely to occur, thus
ignoring the central image which is less deterimental to spotting motion
(change). When confined to a static environment of solitude and
bleakness, the human mind will begin to hallucenate as it seeks a sign
of change - something that says there is still life in the system.

The nature of change is fundamental to physics and philosophy. All
physical laws depend on a description of change (or lack thereof), and
there a wide variety of fields and practices in which the patterns and
structures of dynamical systems can be identified and applied. Dynamical
systems provides a rigorous mathematical framework to describe these
changes, and applying the formalism to other disciplines could yield
curious insights and applications.

Of key significance is the classical relationship between mass,
acceleration, and force - described by Newton's Second Law - and the
relationship between charge, the electric field, and the electrostatic
force - given by Coulomb's Law:

Newton's Second Law: \[F=ma\]

Coulomb's Law: \[F=qE\]

Note both have similar forms, each denoting a force as the product of
two interacting components - a source (mass, \(m\), or charge, \(q\))
and a field (acceleration, \(a\), or electric field \(E\)).
Interestingly enough, one can combine the words ``source'' and ``field''
to create the word ``force''. This also makes for a nice mnemonic. So,
the generalized force can be expressed as:
\(Force =Source\times Field\)\$

Also note that force is measured in Newtons in both cases. Mass is
measured in kilograms and charge is measured in Coulombs. Thus, the
units of acceleration are Newtons per kilogram, or meter per second
squared, and the units of the electric field are Newtons per Coulomb, or
Volts per meter (remember this for later). If the electric field is
considered as analogous to an acceleration field, then the unit for
generalized acceleration is given by: \$Acceleration =
/frac\{Force\}\{Source\} \$

In essence, Sources interact with Fields and experience Forces, and
Forces are the product of an interaction between a Source and a Field.

This generalization can be recognized in Lewin's equation to describe
the behavior of a person in their environment (citation): \(B=f(P,E)\)

which is generally stated as: \(Behavior =Person\times Environment\)

Here, the units of each parameter are more qualitative than
quantitative, though the formalism is easily identifed.

Sources carry intrinsic information such as mass, charge, energy, even
genetics. Fields carry extrinsic information about the environment.
Forces inevitably carry information about both intrinsic and extrinsic
properties of the system. This means that it is not possible to
immediately infer whether a force or outcome is due solely to intrinsic
or extrinsic properties without a controlled experimental variance of
either property to test for such influence.

It is curious to consider the possibility that Fields could be
influenced by Sources themselves. Perhaps the motion of the Source
through the Field had a lasting effect on it, perhaps disspiating Field
Energy or something of the sort - like walking through snow and leaving
behind a trail of deformation. It's worth noting that this may have
something to do with the nature of the Field - as the Field itself has a
Source! So perhaps a Source can have a lasting impact on another.
Consider too that this may relate to the nature of conservative forces
and nonconservative forces and the fields through which they act.*

To explore the nature of this interaction, consider two masses placed a
distance, \(R\), apart from each other. Assuming they are isolated from
the influence of any other source of mass, the force experience by each
mass is the same as described by Newton's Third Law, and the force
between them is described by Newton's Univeral Law of Gravitation:

\(Force = G\frac{mM}{R^2}\)

where \(G\) is the universal gravitational constant, \(M\) is the larger
mass - sometimes called the parent or planet, and \(m\) is the smaller
mass - the child or satellite.

Now consider Coulomb's Law, \(Force = k\frac{qQ}{R^2}\)

where the source terms are now charges \(q\) and \(Q\), and the constant
\(k\) is the Coulomb constant. Note the similarities between these two
equations.

To isolate the effect of the Field on one of the Sources (say, the
satellite), simply rearrange the equation \(Force = G\frac{mM}{R^2}\)
\(\frac{Force}{m} = G\frac{M}{R^2}\) and since the acceleration is the
ratio of force per mass by Newton's Second Law, \(a = G\frac{M}{R^2}\)

This says that the acceleration experienced by \(m\) is caused by \(M\).
So, from the perspective of \(m\), the properties of the acceleration
field are extrinsic and caused by \(M\). The force that \(m\)
experiences in the field also depends on its own intrinsic properties -
specifically mass. The more massive \(m\) becomes, the slower it will
accelerate provided the same field at the same location. For now, it is
assumed that \(m\) cannot influence the field produced by \(M\).

Since Forces are the product of Sources and Fields, they are influenced
by both intrinisic and extrinsic properties.

\section{Potential Energy and
Potential}\label{potential-energy-and-potential}

The configuration of a Source in a Field, and its corresponding Force,
can also be described via its location within the Field. This
description is known as Potential Energy, and can be thought of as a
sort of ``spatial energy''.

Gravitational Potential Energy, \(U = mgh = Fh\)

where the gravitational force by Newton's Second Law is \(F=mg\).

Note that the potential energy is a product of Force and Distance, and
thus has units of \(N\cdot m\) which is a unit of Energy measured in
Joules.

The distance \(h\) is the distance from the field's zero-potential or
``ground''.

A similar expression can be seen in the Electrostatic Potential Energy,
where \(U = qEs = Fs\)

where the electrostatic force is given by \(F=qE\).

If either a mass or a charge (or a mass with charge) is ``elevated'' to
a non-zero potential, there is energy which ``wishes'' to be released.
Upon releasing the mass or charge, or setting it on it's way with some
impetus, it will be driven by a force to follow a trajectory which seeks
to minimize this potential energy.

The force which develops as a result of this potential energy
minimization can be described independently of the source (mass/charge),
more specifically, as the gradient of the potential function.

\(F = -\nabla\phi\)

Define potential as the ratio of potential energy per source (specific
potential energy? energy density?) such that it is the product of Field
and (Reference) Displacement.

For electric potential, commonly known as voltage, \(V=\frac{U}{q}=Es\)

has units of Joules per Coulomb or Volts.

For gravitational potential, not commonly known as gravitational
voltage, \(V=\frac{U}{m}=gh\)

which has units of Joules per kilogram or \(\frac{m^2}{s^2}\). Note
these are the same units as velocity-squared.

\section{Work and Kinetic Energy}\label{work-and-kinetic-energy}

Moving the source of mass/charge through the field requires energy -
specifically at the expense of potential energy. Energy involved in
moving the source from one potential to another is known as Work and is
defined as

\(W = Fdcos\theta\)

where W is measured in Joules {[}J{]}, F is the force measured in
Newtons {[}N{]}, d is the distance in meters {[}m{]}, and θ is the angle
between the applied force and the line of action (direction of motion).

Angles greater than zero indicate suboptimal force influence on the
motion of the source/object. This is more accurately described in the
dot-product relationship of work as

\(W = \int \vec{F}\cdot \mathrm{d}\vec{r}\)

since work is the energy which contributes to change in motion along a
line of action (given by the vector \(\vec{r}\)), this also defines
Kinetic Energy.

\(W= \Delta KE\)

The change in motion is the result of the net work done by the system,
from both conservative forces efficiently transferring energy and
nonconservative forces inefficiently transfering (removing) energy. When
energy is fully conserved, all potential energy is converted into
kinetic energy. However, in real systems, energy is lost in one form or
another to nonconservative forces such as friction.

The work done by these nonservative force is described using the
original defintion,

\(W_{nc} = \int \vec{F}_{nc}\cdot \mathrm{d}\vec{r}\)

While nonconservative forces impact the kinetic energy of the system,
conservative forces and the work they perform depend only on (the loss
of) potential energy.

\(W_{c} = \int \vec{F}_{c}\cdot \mathrm{d}\vec{r} = -\Delta PE\)

The conservation of energy statement fully describes this exchange.

\({KE}_i+{PE}_i +W_{nc}={KE}_f+{PE}_f\)

Thus, for sake of showing consistency in the definitions\ldots{}

The total work is the sum of the conservative and nonconserative work:

\({KE}_i+{PE}_i -{PE}_f+W_{nc}={KE}_f\)

\({KE}_i+-({PE}_f -{PE}_i)+W_{nc}={KE}_f\)

\(-\Delta PE+W_{nc}={KE}_f-{KE}_i\)

\(W_{c}+W_{nc}= \Delta KE = W\)

And the change in energy of the system is zero when energy is conserved.

\(-\Delta PE+W_{nc}={KE}_f-{KE}_i\)

\(W_{nc}=\Delta KE+\Delta PE\)

which says that energy is either added or removed from the system (not
constant). If all of the work is conservative and the nonconservative
work is zero,

\(0=\Delta KE+\Delta PE\)

and energy is conserved (remains constant).

\section{Dynamics}\label{dynamics}

If the energy of a system can be fully described, then its behavior can
be described too.

The energy of a system can be described with three energy modes:

\begin{itemize}
\tightlist
\item
  \textbf{Potential energy}: energy stored in \emph{configuration}
  (position, deformation, separation).
\item
  \textbf{Kinetic energy}: energy stored in \emph{motion}.
\item
  \textbf{Dissipative energy}: energy \emph{lost} to nonconservative
  effects (friction, drag, resistance).
\end{itemize}

Potential energy represents the energy associated with a configuration
in a field, while kinetic energy represents the energy associated with
movement through a field. Dissipative energy is energy lost to
nonconservative effects like friction, drag, and electrical resistance.

Each of these three forms of energy represents a passive role in the
system.

\begin{itemize}
\tightlist
\item
  \textbf{store energy in a potential-like way}\\
\item
  \textbf{store energy in a kinetic-like way}\\
\item
  \textbf{or dissipate energy}
\end{itemize}

\subsection{RLC Circuit}\label{rlc-circuit}

In a classical RLC circuit, the three passive compoenents are the
resistor, inductor, and capacitor. The resistor is the dissipative
element. The inductor is the kinetic storage element (in the magnetic
field, and in the inertia-like behavior of current). The capacitor is
the potential storage element (in the electric field, and in charge
separation).

Before the circuit is closed, the capacitor is uncharged---there is no
voltage across it---and the inductor carries no current. When the switch
closes, current starts to flow. The resistor opposes the flow and
converts electrical power into heat. The inductor opposes changes in the
flow, so it ``pushes back'' against sudden jumps in current. The
capacitor begins accumulating charge, which builds a voltage across it.

As the capacitor charges, the voltage across it rises, leaving less
voltage available to drive current through the loop. The current falls.
In the ideal DC limit, current goes to zero once the capacitor reaches
the source voltage. So the circuit evolves from kinetic activity
(current) into potential storage (charge separation), while the resistor
continuously bleeds energy away.

\subsection{Mass-Spring-Damper}\label{mass-spring-damper}

A similar behavior appears in mechanics with the mass-spring-damper
system.

In a classical mass--spring--damper system, the three passive components
are the damper, mass, and spring. The damper is the dissipative element:
it converts mechanical power into heat through friction-like effects.
The mass is the kinetic storage element: it stores energy in motion and
resists changes in velocity. The spring is the potential storage
element: it stores energy in deformation and ``pushes back'' when
stretched or compressed.

Before anything moves, the spring can be undeformed (no spring force)
and the mass can be at rest (no kinetic energy). When the system is
disturbed---pulled, released, or driven---motion begins. The damper
opposes motion and dissipates energy, producing a resistive force that
grows with velocity. The mass resists rapid changes in velocity, so it
``pushes back'' against sudden acceleration. The spring begins to
stretch or compress, building a restoring force as it stores potential
energy.

As the spring deformation increases, the restoring force grows, reducing
the net force available to keep accelerating the mass. The velocity
tends to peak and then fall as the spring pulls the mass back toward
equilibrium. In the absence of continued driving, the damper steadily
bleeds energy out of the motion, so the oscillations shrink and the
system settles toward rest. So the mechanical system evolves through an
exchange between kinetic energy (mass in motion) and potential energy
(spring deformation), while the damper continuously drains energy away.

This is not just a poetic similarity --- the governing relationships
line up in the same roles:

\begin{itemize}
\tightlist
\item
  The mass resists changes in velocity (a mechanical ``inductor'').
\item
  The spring produces a restoring force proportional to displacement (a
  mechanical ``capacitor,'' in the sense of potential storage).
\item
  The damper produces a resistive force proportional to velocity (a
  mechanical ``resistor'').
\end{itemize}

In other words, the components line up by role:

Resistor (R) ↔ Damper (b): dissipates energy and drains motion/flow into
heat.

Inductor (L) ↔ Mass (m): stores kinetic energy and resists changes in
flow/velocity.

Capacitor (C) ↔ Spring Compliance (1/k): stores potential energy and
builds a restoring ``push'' as it charges/deforms. (Stiffness \(k\) is
the inverse role.)

\subsubsection{Constitutive Laws}\label{constitutive-laws}

The mapping above becomes concrete when you write the element laws in
matching form.

Electrical (one common form): - Resistor: \(V_R = R I\) - Inductor:
\(V_L = L\,\dot I\) - Capacitor: \(I_C = C\,\dot V\)

Mechanical (translation, with \(v=\dot x\)): - Damper: \(F_d = b v\) -
Mass: \(F_m = m\,\dot v \; (= m\ddot x)\) - Spring: \(F_s = kx\)
(equivalently, \(\dot F_s = k v\))

That one line you wrote down --- \(F = m\dot v\) --- is the mechanical
twin of the inductor law \(V = L\dot I\).

Both say: the kinetic storage element resists \emph{changes in the flow
variable}.

\subsection{Restriction--Inertance--Accumulator
(Hydraulics)}\label{restrictioninertanceaccumulator-hydraulics}

A similar behavior appears in hydraulics with a
restriction--inertance--accumulator system.

In a classical hydraulic loop, the three passive components can be taken
as: a restriction (a valve/orifice or narrow section of pipe), an
inertance (a long pipe / moving slug of fluid), and a compliance (an
accumulator with a flexible diaphragm or trapped compressible volume).
The restriction is the dissipative element: it converts mechanical power
into heat through viscous losses and turbulence. The inertance is the
kinetic storage element: the moving fluid has inertia and resists rapid
changes in flow (the same effect behind water hammer). The compliance is
the potential storage element: it stores energy by compressing/expanding
a volume (diaphragm stretch, trapped gas compression, etc.).

Before anything happens, the fluid can be at rest (no kinetic energy),
and the accumulator can be relaxed (no stored pressure energy). When a
pump is engaged or a valve is opened to a pressurized reservoir, flow
begins. The restriction immediately ``pushes back'' against that flow by
demanding a pressure drop to sustain it. The inertance also pushes back,
but differently: it resists sudden changes in flow, so the flow cannot
jump instantly. Meanwhile, the accumulator begins to accept fluid, which
stretches/compresses its compliant element and builds pressure as it
stores potential energy.

As the accumulator pressurizes, its back-pressure rises, leaving less
pressure difference available to keep driving flow through the loop. The
flow rate peaks and then falls as the pressure levels approach
equilibrium. In the steady limit (no oscillatory driving), flow can drop
toward zero even though the system remains pressurized---because the
compliance has saturated or ``charged,'' just like a capacitor.
Throughout the entire process, the restriction steadily bleeds energy
away.

So all three systems share the same basic structure: two elements
exchange energy back and forth (kinetic ↔ potential), while the
dissipative element continuously bleeds energy away. The nouns
change---charge and current versus displacement and velocity---but the
behavior is the same: storage, exchange, decay.

--- AI ---

A useful way to see why these analogies keep working is to start from a
more general lens that already appears in the earlier work:
\textbf{potential is energy per source}. Gravitational potential is
energy per unit mass, and electrical potential is energy per unit
charge.:contentReference{oaicite:0} Once potential is treated as
``energy per source,'' the idea of a \emph{through} quantity becomes
natural: it is a \textbf{source-flow rate} (charge flow, mass flow,
volume flow, etc.). The earlier notes make this explicit by introducing
a gravitational ``Ohm's law'' where \textbf{B is mass flow rate}
(mass-current) and the proportionality constant is a \textbf{flow
resistance} (``Gohs'').:contentReference{oaicite:1}

That same reasoning also produces a clean capacitance-style result: if
\textbf{H} represents a \emph{mass-capacitance}, then the target form is
``source proportional to potential,'' directly mirroring (q =
CV).:contentReference{oaicite:2}:contentReference{oaicite:3}

With that framing, the terminal pairs (voltage/current, force/velocity,
pressure/flow, potential/mass-flow) stop looking arbitrary. They are the
pairs that make power bookkeeping consistent.

\subsubsection{Terminal pairs and the three passive
roles}\label{terminal-pairs-and-the-three-passive-roles}

\begin{longtable}[]{@{}
  >{\raggedright\arraybackslash}p{(\linewidth - 6\tabcolsep) * \real{0.2500}}
  >{\raggedright\arraybackslash}p{(\linewidth - 6\tabcolsep) * \real{0.2500}}
  >{\raggedright\arraybackslash}p{(\linewidth - 6\tabcolsep) * \real{0.2500}}
  >{\raggedright\arraybackslash}p{(\linewidth - 6\tabcolsep) * \real{0.2500}}@{}}
\toprule\noalign{}
\begin{minipage}[b]{\linewidth}\raggedright
Domain / Port
\end{minipage} & \begin{minipage}[b]{\linewidth}\raggedright
Driving variable (effort)
\end{minipage} & \begin{minipage}[b]{\linewidth}\raggedright
Through variable (flow)
\end{minipage} & \begin{minipage}[b]{\linewidth}\raggedright
Power
\end{minipage} \\
\midrule\noalign{}
\endhead
\bottomrule\noalign{}
\endlastfoot
Electrical & (V) & (I) & (P = VI) \\
Mechanical (translation) & (F) & (v) & (P = Fv) \\
Mechanical (rotation) & (\tau) & (\omega) & (P = \tau \omega) \\
Hydraulic & (\Delta p) & (Q) & (P = \Delta p,Q) \\
Gravitational (transport port) & (V\_g) (J/kg) & (B=\dot m) (kg/s) & (P
= V\_g B) \\
\end{longtable}

--- AI ---

At the terminals of a system, there is usually a ``driving difference''
and a resulting ``through'' quantity, corresponding to how the potential
and kinetic elements exchange energy through the system.

In many systems, the ratio of driving difference to through quantity is
set by the dissipative path, and it controls how quickly energy is bled
away relative to how much is stored. This is known as
\textbf{impedance}, and it is defined for each system as the ratio of
``pressure'' to ``flux'':

\begin{itemize}
\tightlist
\item
  Electrical: \(Z = \frac{V}{I}\)
\item
  Mechanical (translation): \(Z = \frac{F}{v}\)
\item
  Mechanical (rotation): \(Z = \frac{\tau}{\omega}\)
\item
  Hydraulic: \(Z = \frac{\Delta p}{Q}\)
\end{itemize}

Note that this is not a ratio of stored energies. It's a ratio of the
\textbf{driving variable} to the \textbf{through variable} at a
port/terminal.

--- AI --- For a purely dissipative path, that ratio is just the
resistive constant. Once storage is present, the ``effective ratio''
becomes dynamic (time / frequency dependent), because part of the
driving difference is temporarily in storage before being returned or
dissipated. --- AI ---

It now becomes natural to describe systems using paired variables like
voltage/current, force/velocity, and pressure/flow --- because these are
the parameter pairs that describe how energy propagates through the
system.

\subsection{Power Accounting}\label{power-accounting}

Because these paired variables describe how energy propagates through
the system, their product conveys a different meaning: \textbf{power},
the rate at which energy is transferred.

\begin{itemize}
\tightlist
\item
  Electrical: \(P = VI\)\\
\item
  Mechanical (translation): \(P = Fv\)\\
\item
  Mechanical (rotation): \(P = \tau \omega\)\\
\item
  Hydraulic: \(P = \Delta p \, Q\)
\end{itemize}

Describing a system at its terminals using \((V,I)\) or \((F,v)\) or
\((\Delta p,Q)\) makes it possible to directly track \textbf{energy
transfer}.

This also clarifies why real sources ``sag under load.'' Real sources
contain internal dissipative effects, so part of the driving difference
is spent internally when the through quantity increases. A battery is
the clean electrical example: it behaves like an EMF in series with an
internal resistance, so the terminal voltage decreases as current
increases:

\(V_{\text{terminal}} = \mathcal{E} - rI\)

A pressurized canister behaves the same way in hydraulics: the reservoir
pressure is the available driving difference, but the outlet pressure
falls as flow increases because pressure is lost across internal
restrictions. Same structure, different nouns.

At this point it becomes reasonable to give these paired variables a
generic name. Bond-graph prior art calls them \textbf{effort} (the
driving difference) and \textbf{flow} (the through quantity), precisely
because their product is power and because the same three passive roles
(potential storage, kinetic storage, dissipation) can be written
consistently in those variables.

This is the bridge to dynamical systems. The moment a system includes
storage --- something that accumulates over time --- its behavior stops
being a purely algebraic ``ratio at the terminal'' and becomes an
evolution law. That is where ODEs (and, when the storage is distributed
in space, PDEs) enter the conversation.

Before discussing differential equations, a table is presented which
collects and organizes the analogy thusfar.

\begin{longtable}[]{@{}
  >{\raggedright\arraybackslash}p{(\linewidth - 10\tabcolsep) * \real{0.1667}}
  >{\raggedright\arraybackslash}p{(\linewidth - 10\tabcolsep) * \real{0.1667}}
  >{\raggedright\arraybackslash}p{(\linewidth - 10\tabcolsep) * \real{0.1667}}
  >{\raggedright\arraybackslash}p{(\linewidth - 10\tabcolsep) * \real{0.1667}}
  >{\raggedright\arraybackslash}p{(\linewidth - 10\tabcolsep) * \real{0.1667}}
  >{\raggedright\arraybackslash}p{(\linewidth - 10\tabcolsep) * \real{0.1667}}@{}}
\toprule\noalign{}
\begin{minipage}[b]{\linewidth}\raggedright
Component / parameter (catalog row)
\end{minipage} & \begin{minipage}[b]{\linewidth}\raggedright
\textbf{Gravitational (potential / mass-flow port)}
\end{minipage} & \begin{minipage}[b]{\linewidth}\raggedright
\textbf{Electrical}
\end{minipage} & \begin{minipage}[b]{\linewidth}\raggedright
\textbf{Mechanical (translation)}
\end{minipage} & \begin{minipage}[b]{\linewidth}\raggedright
\textbf{Mechanical (rotation)}
\end{minipage} & \begin{minipage}[b]{\linewidth}\raggedright
\textbf{Hydraulic}
\end{minipage} \\
\midrule\noalign{}
\endhead
\bottomrule\noalign{}
\endlastfoot
\textbf{Displacement / coordinate} & Height / position \(h\) & Plate
separation / path coordinate \(x\) (contextual) & Position \(x\) & Angle
\(\theta\) & Volume-coordinate / compliance volume \(V\) \\
\textbf{Motion = displacement/time} & Mass-flow speed is contextual;
port flow is \(B=\dot m\) & Carrier drift speed \(u_d\) (contextual) &
Velocity \(v=\dot x\) & Angular speed \(\omega=\dot\theta\) & Volume
flow rate \(Q=\dot V\) \\
\textbf{Source (field-domain ``stuff'')} & Gravitational Mass \(m\) &
Electrical Charge \(q\) & --- & --- & --- \\
\textbf{Field = Force/Source} & Gravitic field \(g\) & Electric field
\(E\) & Acceleration \(a\) & Angular accel. \(\alpha\) &
Pressure-gradient / force-density view (contextual) \\
\textbf{Force = Source·Field} & \(F=mg\) & \(F=qE\) & \(F=ma\) &
\(\tau=J\alpha\) & \(F=pA\) (local relation) \\
\textbf{Work-Energy = Force·Displacement} & \(U=mgh\) & \(U=\int F\,dx\)
(e.g., \(qEs\)) & \(U=\int F\,dx\) & \(U=\int \tau\,d\theta\) &
\(U=\int p\,dV\) \\
\textbf{Potential = Energy/Source} & Gravitational potential
\(V_g=U/m=gh\) & Voltage \(V=U/q\) & (Not typically used as ``per mass''
here) & (Not typically used as ``per mass'' here) & Pressure \(p=U/V\)
(energy density) \\
\textbf{Pressure} & --- & --- & --- & --- & \(p\) (Pa = N/m\(^2\) =
J/m\(^3\)) \\
\textbf{Flux / current} & Mass flow \(B=\dot m\) & Current \(I=\dot q\)
& --- & --- & \(Q=\dot V\) (volume flux) \\
\textbf{Current density / flux density} & Mass flux density (contextual)
& Current density \(J\) (A/m\(^2\)) & --- & --- & Volumetric flux
density \(q_v\) (m/s) (contextual) \\
\textbf{Port ``driving difference'' (bond-graph effort \(e\))} &
Potential \(V_g\) & Voltage \(V\) & Force \(F\) & Torque \(\tau\) &
Pressure difference \(\Delta p\) \\
\textbf{Port ``through quantity'' (bond-graph flow \(f\))} & Mass flow
\(B\) & Current \(I\) & Velocity \(v\) & Angular speed \(\omega\) &
Volume flow \(Q\) \\
\textbf{Power at a port} & \(P=V_g\,B\) & \(P=VI\) & \(P=Fv\) &
\(P=\tau\omega\) & \(P=\Delta p\,Q\) \\
\textbf{Impedance (port ratio)} & \(Z_g=V_g/B\) & \(Z=V/I\) & \(Z=F/v\)
& \(Z=\tau/\omega\) & \(Z=\Delta p/Q\) \\
\textbf{Resistance (R-type coefficient)} & \(R_g\) with \(V_g=R_g B\)
(your ``Gohs''-style slot) & \(R\) with \(V=RI\) & \(b\) with \(F=bv\) &
\(b_\theta\) with \(\tau=b_\theta\omega\) & \(R_h\) with
\(\Delta p=R_h Q\) \\
\textbf{Conductance / mobility (G-type coefficient)} & \(G_g=1/R_g\)
with \(B=G_g V_g\) & \(G=1/R\) with \(I=GV\) & \(\mu=1/b\) with
\(v=\mu F\) & \(\mu_\theta=1/b_\theta\) with \(\omega=\mu_\theta\tau\) &
\(G_h=1/R_h\) with \(Q=G_h\Delta p\) \\
\textbf{Capacitance / compliance (C-type coefficient)} &
Mass-capacitance \(H\) with \(m=H V_g\) and \(B=H\dot V_g\) & \(C\) with
\(q=CV\) and \(I=C\dot V\) & \(C_m=1/k\) with \(x=C_m F\) &
\(C_\theta=1/k_\theta\) with \(\theta=C_\theta\tau\) & \(C_h\) with
\(V=C_h\,\Delta p\) and \(Q=C_h\,\Delta\dot p\) \\
\textbf{Stiffness / elastance (inverse of C)} & \(1/H\) & \(1/C\) &
\(k\) & \(k_\theta\) & \(E_h=1/C_h\) \\
\textbf{Inductance / inertance (I-type coefficient)} & \(L_g\) with
\(\pi_g=L_g B\) and \(V_g=L_g\dot B\) & \(L\) with \(\lambda=LI\) and
\(V=L\dot I\) & Mass \(m\) with \(p=mv\) and \(F=m\dot v\) & Inertia
\(J\) with \(L=J\omega\) and \(\tau=J\dot\omega\) & Inertance \(I_h\)
with \(\Pi=I_h Q\) and \(\Delta p=I_h\dot Q\) \\
\textbf{C-state variable (bond-graph \(q=\int f\,dt\))} & Mass \(m\) &
Charge \(q\) & Displacement \(x\) & Angle \(\theta\) & Volume \(V\) \\
\textbf{I-state variable (bond-graph \(p=\int e\,dt\))} &
Potential-impulse \(\pi_g=\int V_g\,dt\) (formal slot) & Flux linkage
\(\lambda\) & Momentum \(p\) & Angular momentum \(L\) & Pressure impulse
\(\Pi\) \\
\textbf{Unit signature noted in the doc (catalog tags)} &
\textbf{Capacitance:} {[}Source/Potential{]}; \textbf{Conductance:}
{[}Flow/Potential{]}; \textbf{Inductance:} {[}Area/Source{]} & same
tag-structure applied by role & same tag-structure applied by role &
same tag-structure applied by role & same tag-structure applied by
role \\
\end{longtable}

make sure all table parameters are properly dervied/defined

relate to generalized bond-graph systems

analogy between effort/flow and drive/response to lead into dynamical
systems/oscillators/waves




\end{document}
